\chapter{Organización y gestión del proyecto}
\label{cap:organizacion}
\section{Organización}
\label{sec:organizacion}
En esta sección se pretende describir la organización del proyecto que se ha llevado a cabo, desde
la arquitectura más general a los detalles de los componentes individuales y las pruebas
realizadas.

\subsection{Arquitectura del sistema}
El proyecto desarrollado se trata de una aplicación web que permite el registro de usuarios
y ofrece una serie de minijuegos. Estos juegos han sido desarrollados en \texttt{Unity} e
incluidos en una aplicación web de \texttt{React}. La información sobre los usuarios y las
actividades realizadas se almacena en un servidor independiente, al que se accede a través de
una \gls{api}.

La arquitectura del sistema sigue el patrón cliente-servidor, en específico, es una arquitectura
en tres capas. Cada una de las capas tiene una funcionalidad limitada, lo que facilita el
mantenimiento de la aplicación. En este caso concreto nos encontramos con un servidor que
proporciona la funcionalidad de la lógica de negocio, otro que nos provee con el acceso persistente
a datos de la aplicación y un último que sirve y presenta la interfaz gráfica a los usuarios.

Esta diferenciación entre capas se puede ver en la figura \ref{fig:paquetes-clases-diseno}. Cada
capa estaría compuesta por un paquete (a excepción de la interfaz, que se compone tanto del
paquete \enquote{Unity} y \enquote{React}).

\diagrama{
    nombre = {Diagrama de paquetes de clases de diseño},
    archivo = {anexo2/diagramas/clasesDiseñoPaquetes.pdf},
    etiqueta = {fig:paquetes-clases-diseno},
}

\subsection{Diseño de componentes}

\subsection{Proceso de desarrollo}
Puesto que se ha utilizado una metodología ágil para el desarrollo del proyecto, el proceso de
construcción del sistema se ha realizado de manera iterativa e incremental, en vez de por fases muy
específicas. Cada iteración ha supuesto una serie de mejoras y adiciones a distintas partes del
sistema en función de las prioridades y necesidades del proyecto.

Sin embargo, sí que se puede distinguir un \enquote{orden} seguido durante el desarrollo del
sistema. El \gls{backend} fue de las primeras partes del sistema en ser desarrolladas, ya que es
necesario para que la parte del \gls{frontend} funcione correctamente. Por tanto, la construcción
de este último no empezó hasta que el primero estuviese en un estado suficientemente avanzado de
desarrollo. Por otro lado, puesto que el videojuego de \texttt{Unity} es un componente embebido en
la aplicación web con una interfaz con el \gls{frontend} muy específica y limitada, su desarrollo
se realizó prácticamente junto al del \gls{frontend}.

Puesto que \texttt{Jira} permite la asignación de etiquetas a las tareas y que se han utilizado
desde el principio del proyecto para categorizar las tareas según pertenezcan al \gls{backend},
\gls{frontend}, \texttt{Unity} o a la documentación, se ha creado un gráfico que muestra la
cantidad de puntos de historia completados (acumulados) en cada sprint, separados por
categorías (figura \ref{fig:sprintstags}). Como se puede ver, la mayoría de los puntos de historia
completados pertenecen al videojuego, seguido de la documentación, el \gls{frontend} y, por último,
el \gls{backend}.

\diagrama{
    nombre = {Puntos de historia completados por categoría},
    archivo = {memoria/fotos/sprintstags.pdf},
    etiqueta = {fig:sprintstags},
    width = 1\textwidth,
}


\subsection{Pruebas y verificación}
Durante el desarrollo del proyecto era necesario evaluar el correcto funcionamiento de la aplicación,
así como la precisión de la detección de posturas. Aunque para el segundo de estos aspectos el
proceso de validación consistió en probar la aplicación múltiples veces durante el desarrollo de
la misma, para el primero se definieron una serie de pruebas unitarias para garantizar el correcto
funcionamiento de la \gls{api}. Las pruebas unitarias se comenzaron a definir en los primeros sprints
del desarrollo, una vez se creó la estructura de la \gls{api} y se implementaron las primeras
funciones. Estas pruebas, que se ejecutan de manera casi automática, permiten verificar que cada nueva
versión del código no introduce errores en las funcionalidades ya implementadas.


\section{Gestión del proyecto}
\label{sec:gestion}
En esta sección se desarrollarán tanto los recursos utilizados durante el desarrollo del proyecto y
un resumen de la planificación temporal detallada en el \tercerAnexo.

\subsection{Recursos utilizados}
% Hardware, software y recursos personales
Podemos agrupar los recursos utilizados en tres categorías: recursos de hardware, software y
recursos personales.

Los recursos de hardware se refieren a los equipos físicos utilzados para construir el sistema. En
este caso se trata de mi ordenador portátil (HP Victus 16-e0xxx\footnote{https://support.hp.com/ar-es/drivers/victus-by-hp-16.1-inch-gaming-laptop-pc-16-e0000/model/2100371558?sku=6X738U8})
y mi teléfono móvil (Samsung Galaxy A35 5G\footnote{https://www.samsung.com/levant/smartphones/galaxy-a/galaxy-a35-5g-awesome-lemon-256gb-sm-a356ezyvmea/}).
Ambos dispositivos tienen cámaras integradas, lo que hace posible utilizarlos para probar el
producto.

Entre los recursos de software se encuentran aquellos programas y herramientas utilizados durante el
desarrollo del proyecto. No se mencionarán aquí bibliotecas o \glspl{framework} utilizados, sino herramientas
más genéricas. Para una descripción de todas las herramientas del código ver el apartado \ref{sec:herramientas}

Para la escritura del código se han utilizado los editores \tech{Code}{https://code.visualstudio.com/}
y \tech{NeoVim}{https://neovim.io/}. Para probar la aplicación creada se han usado los navegadores web
\tech{Firefox}{https://www.mozilla.org/en-US/firefox/new/} y \tech{Chromium}
{https://www.chromium.org/chromium-projects/}. Para el desarrollo del videojuego se ha utilizado el motor
Unity mencionado anteriormente. Por último, como sistema de control de versiones se han utilizado tanto
\tech{Git}{https://git-scm.com/} como \tech{GitHub}{https://github.com/}.

Por último, los recursos humanos utilizados me incluirían a mí y a los tutores del proyecto. Todas las
labores de desarrollo, diseño y pruebas han sido realizadas por mí, mientras que los tutores han
supervisado el proyecto y han proporcionado orientación y apoyo en la toma de decisiones.

\subsection{Planificación temporal}
La planificación del proyecto se ha hecho siguiendo la metodología de desarrollo ágil Scrum. De manera
resumida, Scrum organiza el trabajo en iteraciones cortas denominadas \textit{sprints}, que suelen durar
en torno a dos semanas. Cada sprint comienza con una reunión de planificación en la que se definen las
tareas a realizar durante este. Durante el desarrollo del sprint, se realizan reuniones diarias de
seguimiento para revisar el progreso del equipo de trabajo. Al final del sprint, se realizan varias
reuniones de revisión y retrospectiva para evaluar el trabajo realizado y planificar el siguiente sprint
en función de los resultados obtenidos.

Para este proyecto, se han llevado a cabo un total de 22 sprints, cada con una duración de una semana.
En cada uno de ellos se han definido una serie de tareas a realizar, que se habrán o no completado al final
del sprint. En el \tercerAnexo se puede ver un desglose de las tareas planificadas y realizadas en cada
uno de los sprints.

Para dar una visión general de los sprints realizados, se ha creado un diagrama que muestra los puntos
de historia completados (figura \ref{fig:sprints}). Los puntos de historia son una medida de la
complejidad de las tareas realizadas y es la medida principal utilizada para planificar el trabajo de
cada sprint (tal vez detrás de la prioridad de cada tarea). En el gráfico se puede observar que se comenzó
con un ritmo más lento, pero a medida que se fueron completando tareas y se fue adquiriendo experiencia,
el ritmo de trabajo aumentó. También se puede observar la época de exámenes (sprints 16 y 17) en la que
se redujo la cantidad de puntos completados y como a partir del sprint 18 se volvió al proyecto con un
ritmo muchas veces mayor que al principio, puesto que no la totalidad del tiempo se dedicaba al proyecto.

\diagrama{
    nombre = {Todos los sprints realizados},
    archivo = {memoria/fotos/sprints.pdf},
    etiqueta = {fig:sprints},
    width = 1\textwidth,
}

